% Provisional evidence manuscript. Not the official IEICE class or an accepted contribution.
\documentclass[uplatex,a4paper,twocolumn,10pt]{jsarticle}
\usepackage[margin=20mm]{geometry}
\usepackage[dvipdfmx]{graphicx}
\usepackage{amsmath}
% Generated from audited reports; edit the builder, not these tables.
\newcommand{\EvidencePThirtyThreeTotals}{%
\begin{tabular}{lrrr}\hline 指標 & 全表示 & 省略 & 一覧付\\\hline
成功 & 2/4 & 2/4 & 2/4\\
初回入力 & 48,485 & 21,893 & 25,092\\
継続総token & 242,640 & 377,272 & 327,269\\
AI呼出し & 18 & 39 & 29\\
継続API & 162 & 133 & 150\\
実験用復元API & 51 & 51 & 51\\
処理入力token & 57,548 & 43,326 & 50,702\\
出力token & 8,752 & 11,138 & 11,686\\
継続時間(s) & 1,025.8 & 1,043.4 & 1,136.3\\
\hline\end{tabular}}
\newcommand{\EvidencePThirtyThreeTasks}{%
\begin{tabular}{llrr}\hline 課題 & 条件・成否 & token & AI\\\hline
写真整理 & F・成功 & 66,851 & 5\\
 & M・成功 & 63,575 & 6\\
 & I・成功 & 87,160 & 8\\
共有・依頼 & F・未達 & 110,922 & 9\\
 & M・未達 & 110,906 & 12\\
 & I・未達 & 109,531 & 9\\
返金 & F・成功 & 31,680 & 2\\
 & M・成功 & 98,106 & 12\\
 & I・成功 & 27,808 & 4\\
曲一覧書出 & F・未達 & 33,187 & 2\\
 & M・未達 & 104,685 & 9\\
 & I・未達 & 102,770 & 8\\
\hline\end{tabular}}
\newcommand{\EvidencePFortySevenTasks}{%
\begin{tabular}{llrr}\hline 課題 & 条件・成否 & token & AI\\\hline
写真整理 & F・成功 & 77,478 & 6\\
 & N・未達 & 103,158 & 9\\
 & S・未達 & 113,762 & 10\\
共有・依頼 & F・未達 & 115,851 & 9\\
 & N・未達 & 106,952 & 9\\
 & S・未達 & 106,092 & 9\\
返金 & F・成功 & 32,145 & 2\\
 & N・成功 & 63,118 & 4\\
 & S・成功 & 47,185 & 3\\
曲一覧書出 & F・未評価 & 68,964$+?$ & 5\\
 & N・未実行 & --- & ---\\
 & S・未実行 & --- & ---\\
\hline\end{tabular}}

\title{持続する実行状態と入力省略の費用\\AIハーネスの開発比較から投稿の問いを絞る}
\author{検討用原稿・著者情報未設定}
\date{2026年10月22日}
\begin{document}
\raggedbottom
\maketitle
\noindent\textbf{原稿の位置づけ．}
本稿は監査済みの開発結果をつなぐ検討用原稿である．
新しい制御方式の有効性，独立した研究上の貢献，公式書式への適合は未確定であり，
このまま投稿できる完成稿ではない．

\noindent\textbf{あらまし案．}
実行状態を残す出力省略を４開発課題で比較した．初回入力の短縮に対して継続総トークンが増える例を観測し，最終成否，停止理由，修復の適用範囲を分けて示す．

\section{目的と投稿へ向けた問い}
AIオーケストレーションでは，モデルへ渡す情報だけでなく，
情報取得，生成，実行，検証を含む仕事全体を制御する必要がある．
本研究の最終目標は，必要な情報と処理を残し，要求を満たす最終成果の品質を維持しながら，
不要な取得・共有・AI呼出し・実行・検証を減らすハーネス制御である．
短い入力や少ないAPI回数は，この目標の部分的な指標にすぎない．

今回の検討対象は単一作業役の継続である．
同じ途中状態から表示だけを変えたとき，初回の短縮が終了までの費用へ届くか，
届かない場合にどの費用と処理が残るかを調べる．
複数担当間の通信削減や再編の効果を測ったものではない．
また，過去に失敗した条件だけを除くことや，成功まで上限を増やすことは行わない．

\newpage
\section{関連研究と未確定の差}
文献\cite{beyond}は圧縮の方式・時点・量を分け，総費用や時間，課題別成否を調べる．
文献\cite{compression}は状態の再取得と実行を分け，保持内容の介入も行う．
したがって，圧縮後の費用増や同数成功だけでは不十分という一般論は，
本研究の新しい発見として扱わない．

コードの状態を保持することも既存の機能である\cite{codeact,cave}．
本稿で区別するのは，会話中の表示，Python内で利用可能な値，
外部環境への作用という三つの層である．
ただし層を分けただけでは独自の方式や知見は成立しない．
投稿の中心にできるかは，既存の説明では足りない条件や効果を，
適切な対照によって新たに確かめられるかに依存する．

\section{費用と保持状態の定義}
途中状態を，会話履歴$H$，実行名前空間$P$，操作先環境$E$，残予算$B$で記す．
表示省略は$H$の一部を変える．$P$に値が残っていることと，
モデルがその値の存在・用途を把握して使えることは同じではない．
一方，使われた記録がないという理由だけで，その情報を不要とは判定しない．

継続の論理トークンを
\begin{equation}
 T=\sum_{j=1}^{n}(I_j+O_j)
\end{equation}
とする．キャッシュされた入力も$I_j$へ含む．
実処理入力，出力，API，時間，制御・準備費用を別に記録する．
共通前半は仕事全体の費用へ必要だが，分岐後の差とは分ける．
実験のための状態復元を，運用時のAI生成費用へ混ぜない．

\clearpage
\section{同じ途中状態からの開発比較}
AppWorld\cite{appworld}の既知４課題について，持続Python基準の第８応答後を共通の分岐点とした．
課題は写真整理，共有・依頼，返金，曲一覧の書出である．
元の要求・返却，保存された境界，環境とプログラムの対応を照合し，
状態復元費用を別記した．Pythonプロセスのヒープや乱数状態を丸ごと複製したという主張ではない．
保留している12シナリオの要求文は，方式選定のために閲覧していない．

\subsection{出力省略の比較（P33）}
全表示F，古い成功出力の省略M，省略に変数・API一覧を加えるIを比較した．
省略は分岐点で一度だけ行い，最新の完全な応答組，操作コード，
エラー，停止・完了の記録を保つ．以後は通常どおり追記する．
原記録は全条件で読み返せる．一覧は名前，基本型，長さ，公開API名を提示し，
引数値の要約や未来の関連性判定を行わない．

同じモデルMiniMax-M3 IQ4\_XSと継続seedを用い，
課題順に３条件の実行順を回転させた．
継続上限は16生成，120,000 token，480秒で，
API枠は共通前半を含めた元課題の上限とした．
全12条件，各１試行の開発比較である．
この実験は，後に導入した初回キャッシュ不使用の統一より前であり，
同じseedだけで生成経路の再現性を保証しない．

\begin{center}
\begin{minipage}{\columnwidth}\centering
\small 表1\quad P33の全４課題合計\par\smallskip
\EvidencePThirtyThreeTotals
\end{minipage}
\end{center}

\newpage
\subsection{仕様表示と拒否返信の比較（P47）}
P33後の診断を受け，保存された表示処理とAPI返却・出力全体の由来が一致する，
古い仕様部分だけを選択表示の対象にした．
全表示＋詳細契約F，選択表示＋引数名N，選択表示＋詳細契約Sの３条件を設けた．
詳細契約を全表示側にも与え，片側だけの修復情報の追加を省略効果にしない．

返信対象は既公開APIで未知引数のため送信前に拒否された要求である．
元の契約を改変せず返し，自動の引数修正や再送は行わない．
未公開API，送信済み，作用不明の要求は対象外とする．
後続の生成・実行・修復も費用に含む．
残予算通知，残量に応じた出力枠，初回キャッシュ不使用を全条件へそろえた．
上限はP33と同じだが，制御と表示規則が異なるため，
P33との結果の差を単一処理の因果効果へ変換しない．

\subsection{品質と欠測の扱い}
最終環境に対する独立評価で課題ごとの成否を記録する．
開発上の採用基準は，対照の成功を失わず，既知の継続総tokenを減らすこととした．
少数で同数成功となっても統計的な非劣性ではない．
使用量不明はゼロで補わず，事前規則で停止した場合は未評価と未実行を残す．
安い未完了を効率のよい成功とは数えない．

\clearpage
\section{全予定条件の結果}
\subsection{P33：初回入力と終了までの費用}
全12条件を評価し，使用量不明は０だった．
初回入力は短くなったが，継続総tokenはMで55.5\%，Iで34.9\%増えた．
全条件で同じ２課題が成功し，成功課題の入替えはなかった．
一方，書出課題のFはAPI上限，MとIはtoken予約上限で終了しており，
合計値だけで全表示が優れると一般化できない．

\begin{center}
\begin{minipage}{\columnwidth}\centering
\small 表2\quad P33の全12条件\par\smallskip
\EvidencePThirtyThreeTasks
\end{minipage}
\end{center}

返金ではFの２ APIに対し，Mは仕様説明・相手の探索等を含む８ API，Iは４ APIだった．
写真整理ではMはFより低token，Iは高tokenとなった．
一つの成功例や各課題の最良条件を，運用時に選べる制御規則とはみなさない．
全体のAPI減少にも，書出の未完了経路の違いが含まれる．

\newpage
\subsection{P47：契約返信を加えた比較}
９条件を評価した．書出Fの最後の生成が時間上限で使用量不明となり，
事前規則に従って残る２条件を実行しなかった．
全66生成の既知使用量834,705 tokenに不明１呼出しが残る．
継続325 APIと実験用復元117 APIを区別した．
この結果を全４課題の完全な削減率へ変換しない．

\begin{center}
\begin{minipage}{\columnwidth}\centering
\small 表3\quad P47の全12予定条件\par\smallskip
\EvidencePFortySevenTasks
\end{minipage}
\end{center}

写真整理はFで成功し，NとSで未達だった．
返金は３条件で成功したが，SもFより15,040 token多かった．
全体で付加した６返信は次のモデル入力に存在し，
次のプログラムで同じAPIが送信された．
ただし写真Sでは適用対象の拒否がなく，契約返信は０件だった．
返信の存在，同じAPIの後続送信，最終成功を区別する必要がある．

\clearpage
\section{制御規則へつなげる際の限界}
\subsection{費用の内訳は独立した介入効果ではない}
初回入力を$I$，$G=\sum_j(I_j-I)$とすると，$T=nI+G+\sum_jO_j$である．
初回入力の差，生成回数の差，後続入力の増減，出力の差へ恒等的に分けられるが，
各項は独立した因果効果でも，そのまま回収できる節約量でもない．
P33のAPI減少を再取得の解消と解釈することも，
P47の詳細返信後の生成を全て削減可能な修復費用とすることもできない．

\subsection{小さい修復入力が適用できる範囲}
P47の10形式拒否を保存記録で点検した．
値と既存引数を保持する部品では，空契約の通常正規化が３件，
型上の完全な名前対応が１件，完全対応なしが２件，未公開APIが４件だった．
介入前のP24全６とP29全４実行でも点検したところ，
前者に形式拒否はなく，後者の６拒否に型上の名前対応２件があった．
両群の課題は異なり，インターフェース間の因果比較ではない．

追加の２件はログインAPIで，型上の行先が各１つだった．
したがって全拒否を別LLMへ送る前に，通常の正規化や決定的な対応付けを対照へ含める必要がある．
唯一の型候補でも意味的な正解とは限らない．
契約・拒否要求・エラーからAPIを限定して修復する通常方式\cite{schema}との差を明確にし，
辞退，通常継続への復帰，合流後の最終品質まで測る必要がある．
これらの点検は生成・実行を伴わず，節約実績を示していない．

\newpage
\subsection{投稿の中心へ進むために必要な比較}
本稿の開発結果は，実行状態を残した入力省略にも費用と品質の確認が必要であることを示す．
しかし，既に知られた圧縮の注意点を超える貢献が確立したとは言えない．
次の比較には，(i)通常方式でも残る具体的な費用，
(ii)実行時に使える情報だけで必要処理を保つ判断規則，
(iii)同じ情報と予算を持つ強い対照，
(iv)規則が不要・有害となる条件をそろえる必要がある．

特に，状態が失われたための再取得と，保持された状態の利用方法を決める生成，
新たに必要な情報取得を混同しない．
確認群は機構・比較・予算・評価を固定してから使用する．
12シナリオで悪化０件を観測しても，一般的な品質維持の証明にはならない．
単一モデル・既知４課題・各１回，P33の初回推論状態，P47の欠測，
自作接続への依存は主張を制限する．

\section{現在の結論}
本検討では，品質を維持した総費用削減として採用できる出力省略方式は得られなかった．
初回入力，総token，API，時間，最終成否，停止理由を分けた全件記録を，
次の具体的な制御規則と対照を選ぶ根拠にする．
投稿の主張とあらましを確定する前に，独自の追加価値が残る比較が必要である．

\begin{thebibliography}{9}\small
\bibitem{beyond} R. Satish et al., ``Beyond Token Savings: A Systematic Study of Context Compression in LLM Agents,'' arXiv:2609.32961v1, 2026.
\bibitem{compression} S. Liu, ``What Does Context Compression Cost an Agent? Interaction Costs Unrevealed by Task-Completion Metrics,'' arXiv:2608.16370v1, 2026.
\bibitem{codeact} X. Wang et al., ``Executable Code Actions Elicit Better LLM Agents,'' ICML, PMLR 235, pp.50208--50232, 2024.
\bibitem{cave} M. Ran et al., ``CaveAgent: Transforming LLMs into Stateful Runtime Operators,'' arXiv:2601.01569v4, 2026.
\bibitem{schema} A. Chattopadhyay et al., ``Robust Function Calling for LLM Agents Under Interface Drift,'' arXiv:2610.23342v1, 2026.
\bibitem{appworld} H. Trivedi et al., ``AppWorld: A Controllable World of Apps and People for Benchmarking Interactive Coding Agents,'' ACL, pp.16022--16076, 2024.
\end{thebibliography}
\end{document}
